\subsection{分圆扩张}

设 \(p\) 为 \(\mathbf{K}\) 的特征指数，且设 \(n \geqslant 1\) 为与 \(p\) 互素的整数；所谓级数为 \(n\) 的分圆扩张（ \(\mathbf{K}\) 上的），我们指的是任一分裂扩张 \(\mathbf{E}\) （即多项式 \(\mathbf{X}^n - 1\) 在 \(\mathbf{K}\) 上的分裂扩张）（V，第21页）。由于该多项式是可分的， \(\mathbf{E}\) 是 \(\mathbf{K}\) 的伽罗瓦扩张，且次数有限（V，第57页）。存在一个本原的 \(n\) 次单位根（在 \(\mathbf{E}\) 中）；若 \(\zeta\) 是这样一个根，则每个 \(n\) 次单位根都是 \(\zeta\) 的幂，因此 \(\mathbf{E} = \mathbf{K}(\zeta)\) 。

在本号其余部分中，我们将选取 K 的一个可分闭包 \(K_{s}\) 。对每个与 p 互素的 \(n \geqslant 1\) ，群 \(\mu_{n}(K_{s})\) 是 n 阶循环群，且域

\[
\mathrm{R} _ {n} (\mathrm{K}) = \mathrm{K} \left(\mu_ {n} \left(\mathrm{K} _ {s}\right)\right)
\]

是级数为 \(n\) 的、 \(K\) 上的一个分圆扩张。我们可以把 \(\mu_n(K_s)\) 看作环 \(Z/nZ\) 上的秩为1的自由模，且每个元素 \(\sigma\) （属于 \(\text{Gal}(K_s/K)\) ）诱导 \(\mu_n(K_s)\) 的一个自同构；因此存在同态 \(\chi_n: \text{Gal}(K_s/K) \to (Z/nZ)^*\) ，它由公式 \(u(\zeta) = \zeta^j\) 刻画：对每个 \(n\) 次单位根 \(\zeta\) （在 \(K_s\) 中）、每个 \(u\) （在 \(\text{Gal}(K_s/K)\) 中）以及每个整数 \(j\) （属于剩余类 \(\chi_n(u) \mod n\) ）。由于 \(R_n(K) = K(\mu_n(K_s))\) ， \(\chi_n\) 的核是 \(\text{Gal}(K_s/R_n(K))\) （ \(\text{Gal}(K_s/K)\) 的一个子群）；因此我们有 \(\chi_n = \varphi_n \circ \psi_n\) ，其中 \(\psi_n\) 是 \(\text{Gal}(K_s/K)\) 到 \(\text{Gal}(R_n(K)/K)\) 上的限制同态，而 \(\varphi_n\) 是 \(\text{Gal}(R_n(K)/K)\) 到 \((Z/nZ)^*\) 中的一个单同态。特别地，我们有以下结果：

命题5. --- 对每个整数 \(n \geqslant 1\) （与 \(p\) 互素），扩张 \(\mathbf{R}_n(\mathbf{K})\) 是 \(\mathbf{K}\) 的一个次数有限的阿贝尔扩张，其伽罗瓦群同构于 \((\mathbf{Z}/n\mathbf{Z})^*\) 的一个子群，且其次数整除 \(\varphi(n)\) （即 \((\mathbf{Z}/n\mathbf{Z})^*\) 的阶）。

设 \(\bar{\mathbf{Q}}\) 为有理数域 \(\mathbf{Q}\) 的一个代数闭包，并令 \(n\) 为一个整数。则分圆多项式 \(\Phi_n\) （级数为 \(n\) 的）由下式定义

\[
\Phi_ {n} (\mathbf {X}) = \prod_ {\zeta \in S _ {n}} \left(\mathbf {X} - \zeta\right),\tag{5}
\]

其中 \(S_{n}\) 是 \(\bar{Q}\) 中 n 次本原单位根组成的集合。多项式 \(\Phi_{n}\) 的次数为 \(\varphi(n)\) （命题4）。显然 \(\Phi_{n}(X)\) 在 \(\bar{Q}\) 的所有自同构下不变，因此属于 Q{[}X{]}。由于每个元素 \(\zeta\) （ \(S_{n}\) 的）都是多项式 \(X^{n}-1\) 的根，多项式 \(\Phi_{n}(X)\) 整除 \(X^{n}-1\) ，且以下引理表明 \(\Phi_{n}(X)\) 是一个整系数的首一多项式。

引理2. --- 设 \(f, g\) 和 \(h\) 是 \(\mathbf{Q}[\mathbf{X}]\) 中的首一多项式，使得 \(f = gh\) 。若 \(f\) 具有整数系数，则 \(g\) 和 \(h\) 亦然。

设 \(a\) （相应地 \(b\) ）是整数 \(\alpha \geqslant 1\) （相应地 \(\beta \geqslant 1\) ）中最小的、使得 \(\alpha g\) （相应地 \(\beta h\) ）具有整数系数的那一个；我们令 \(g' = ag\) 和 \(h' = bh\) ，并用反证法证明 \(a = b = 1\) 。若非如此，则存在素数 \(p\) 整除 \(ab\) 。若 \(u \in \mathbf{Z}[\mathbf{X}]\) ，我们用 \(\bar{u}\) 表示系数在域 \(\mathbf{F}_p\) 中的、由（模 \(p\) ）约化 \(u\) 的系数而得到的多项式。我们有 \(g'h' = abf\) ，从而 \(\bar{g}'\bar{h}' = 0\) ；由于环 \(\mathbf{F}_p[\mathbf{X}]\) 是整环（IV，第9页，命题8），我们有 \(\bar{g}' = 0\) 或 \(\bar{h}' = 0\) 。换言之， \(p\) 整除 \(g'\) 的所有系数或 \(h'\) 的所有系数，这与假设矛盾。

我们有关系

\begin{enumerate}
\def\labelenumi{(\arabic{enumi})}
\setcounter{enumi}{5}
\tightlist
\item
\end{enumerate}

\[
\mathbf {X} ^ {n} - 1 = \prod_ {d \mid n} \Phi_ {d} (\mathbf {X}).
\]

因为我们有 \(\mathbf{X}^n - 1 = \prod_{\zeta \in \mu_n(\mathbf{Q})} (\mathbf{X} - \zeta)\) ，且集合 \(S_d\) （其中 \(d\) 整除 \(n\) ）构成 \(\mu_n(\mathbf{Q})\) 的一个划分。

公式（6）确定了 \(\Phi_n(\mathbf{X})\) ，只要我们已知 \(\Phi_d(\mathbf{X})\) （对所有满足 \(d < n\) 的、 \(n\) 的因数 d）；由于 \(\Phi_1(\mathbf{X}) = \mathbf{X} - 1\) ，我们因此有一个计算 \(\Phi_n\) 的递归过程。例如，对素数 \(p\) 我们有

\[
\mathbf {X} ^ {p} - 1 = (\mathbf {X} - 1) \Phi_ {p} (\mathbf {X}),
\]

由此

\begin{enumerate}
\def\labelenumi{(\arabic{enumi})}
\setcounter{enumi}{6}
\tightlist
\item
\end{enumerate}

\[
\Phi_ {p} (\mathbf {X}) = \mathbf {X} ^ {p - 1} + \mathbf {X} ^ {p - 2} + \dots + \mathbf {X} + 1,
\]

和

\[
\Phi_ {p ^ {r + 1}} (\mathbf {X}) = \Phi_ {p} \left(\mathbf {X} ^ {p ^ {r}}\right) \quad \text { for } \quad r \geqslant 0.
\]

我们列出多项式 \(\Phi_n(\mathbf{X})\) （对 \(1 \leqslant n \leqslant 12\) ）的值：

\[
\Phi_ {1} (\mathbf {X}) = \mathbf {X} - 1
\]

\[
\Phi_ {2} (\mathbf {X}) = \mathbf {X} + 1
\]

\[
\Phi_ {3} (\mathbf {X}) = \mathbf {X} ^ {2} + \mathbf {X} + 1
\]

\[
\Phi_ {4} (\mathbf {X}) = \mathbf {X} ^ {2} + 1
\]

\[
\Phi_ {5} (\mathbf {X}) = \mathbf {X} ^ {4} + \mathbf {X} ^ {3} + \mathbf {X} ^ {2} + \mathbf {X} + 1
\]

\[
\begin{array}{l} \Phi_ {6} (\mathbf {X}) = \mathbf {X} ^ {2} - \mathbf {X} + 1 \\ \Phi_ {7} (\mathbf {X}) = \mathbf {X} ^ {6} + \mathbf {X} ^ {5} + \mathbf {X} ^ {4} + \mathbf {X} ^ {3} + \mathbf {X} ^ {2} + \mathbf {X} + 1 \\ \Phi_ {8} (\mathbf {X}) = \mathbf {X} ^ {4} + 1 \\ \Phi_ {9} (\mathbf {X}) = \mathbf {X} ^ {6} + \mathbf {X} ^ {3} + 1 \\ \Phi_ {1 0} (\mathbf {X}) = \mathbf {X} ^ {4} - \mathbf {X} ^ {3} + \mathbf {X} ^ {2} - \mathbf {X} + 1 \\ \Phi_ {1 1} (\mathbf {X}) = \mathbf {X} ^ {1 0} + \mathbf {X} ^ {9} + \mathbf {X} ^ {8} + \mathbf {X} ^ {7} + \mathbf {X} ^ {6} + \mathbf {X} ^ {5} + \mathbf {X} ^ {4} + \mathbf {X} ^ {3} + \mathbf {X} ^ {2} + \mathbf {X} + 1 \\ \Phi_ {1 2} (\mathbf {X}) = \mathbf {X} ^ {4} - \mathbf {X} ^ {2} + 1 \end{array}
\]

\(\Phi_1, \Phi_2, \Phi_3, \Phi_4, \Phi_5, \Phi_7, \Phi_8, \Phi_9\) 和 \(\Phi_{11}\) 的值直接由（7）得出；现在我们有 \(\Phi_1\Phi_2\Phi_3\Phi_6 = X^6 - 1\) 和 \(\Phi_1\Phi_2\Phi_3\Phi_4\Phi_6\Phi_{12} = X^{12} - 1\) ，从而 \(\Phi_4\Phi_{12} = \frac{X^{12} - 1}{X^6 - 1} = X^6 + 1\) ，最终 \(\Phi_{12} = \frac{X^6 + 1}{X^2 + 1} = X^4 - X^2 + 1\) 。情形 \(n = 6\) 和 \(n = 10\) 可类似处理。

注记. --- * 对每个整数 \(n>0\) ，函数 \(\mu(n)\) 定义如下：如果 \(n\) 被一个素数的平方整除，则令 \(\mu(n)=0\) ，否则 \(\mu(n)=(-1)^{h}\) ，若 \(n\) 是 \(h\) 个互不相同的素数之积（“默比乌斯函数”）。可以证明

\[
\Phi_ {n} (\mathbf {X}) = \prod_ {d \mid n} \left(\mathbf {X} ^ {n / d} - 1\right) ^ {\mu (d)},\tag{8}
\]

或更明确地

\[
\Phi_ {n} (\mathbf {X}) = \prod_ {p _ {1} <   \dots <   p _ {h}} \left(\mathbf {X} ^ {n / p _ {1} \dots p _ {h}} - 1\right) ^ {(- 1) ^ {h}}\tag{9}
\]

其中 \((p_{1},\ldots,p_{h})\) 遍历 n 的所有严格递增的素因子序列（参见 Lie, II, 第207页，习题1）。*

